\titledquestion{Multiple Choices}

Each question has \textbf{one or more} correct answer(s). Select all the correct answer(s). For each question, you will get 0 points if you select one or more wrong answers, but you will get 1 point if you select a non-empty subset of the correct answers.

Write your answers in the following table.

% Enter selected choices in the blank table below.

\begin{table}[h]
\centering
\begin{tabular}{|p{1.5cm}|p{1.5cm}|p{1.5cm}|p{1.5cm}|}
        \hline
        (a) & (b) & (c) & (d)  \\
        \hline
          B  &  AB   &  ABC   &     AB     \\[0.7cm]
        \hline
        (e) & (f) & (g) & (h) \\
        \hline
          BC  &  AB   &  ABC   &     AC    \\[0.7cm]
        \hline
    \end{tabular}
\end{table}

\textbf{Note that in this problem, heaps are binary heaps, array-based heap are implemented based on complete binary tree.}

\begin{parts}

\part[3] Which of the following statements about the heap is/are true?
    \begin{choices}
        \choice If a binary tree is a max-heap, then the pre-order traversal of this tree is a descending sequence.
        \choice  For any node A in a heap, the subtree of the tree with root A is still a heap.
        \choice We can sort an array in linear time through heap sort since we can build a heap from an array in O(n) time through Floyd’s method.
        \choice If we pop the minimum element from a binary min-heap and insert the same element back, the heap structure always remains unchanged.
    \end{choices}

\part[3] Which of the following statements about array-based max-heaps is/are true?
    \begin{choices}
        \choice In a zero-indexed heap array, the children of node $i$ are stored at indices $2i+1$ and $2i+2$, whenever these indices exist.
        \choice The array $\langle 50,30,40,10,20,35,37\rangle$ can represent a valid zero-indexed max-heap.
        \choice Replacing the root of a max-heap with any smaller value always preserves the max-heap property.
        \choice Every leaf in a max-heap has a value smaller than every internal node.
    \end{choices}

\part[3] Which of the following statements about heap operations is/are true?
    \begin{choices}
        \choice Building a heap from an arbitrary array using Floyd's bottom-up method takes $O(n)$ time.
        \choice Performing $n$ extract-min operations on a binary min-heap takes $O(n\log n)$ time in the worst case.
        \choice Increasing the key of a node in a max-heap can be performed in $O(\log n)$ time if the node position is known.
        \choice Finding the minimum element in a max-heap always takes $O(1)$ time.
    \end{choices}

\newpage
\part[3] Which of the following statements about binary search trees are true?
    \begin{choices}
        \choice The in-order traversal of a BST containing distinct keys is strictly increasing.
        \choice The pre-order traversal of a BST containing distinct keys uniquely determines the BST.
        \choice Searching for a key in an arbitrary BST with $n$ nodes always takes $O(\log n)$ time.
        \choice The in-order traversal of every BST is a decreasing sequence.
    \end{choices}

\part[3] Which of the following statements about BST search and deletion is/are true?
    \begin{choices}
        \choice Every BST with $n$ nodes has height $O(\log n)$.
        \choice If a node has a right subtree, its in-order successor is the leftmost node in that right subtree.
        \choice If a node has no right subtree, its in-order successor, if it exists, is the lowest ancestor for which the node lies in the left subtree.
        \choice The root of every BST is the median key.
    \end{choices}

\part[3] Which of the following statements about Huffman coding is/are true?
    \begin{choices}
        \choice Huffman's algorithm repeatedly combines the two symbols or subtrees with the smallest frequencies.
        \choice A Huffman code minimizes the weighted average codeword length among all binary prefix codes for the given frequencies.
        \choice Huffman coding always assigns codewords of the same length to all symbols.
        \choice The Huffman coding tree generated for a set of frequencies is always unique.
    \end{choices}

\part[3] Which of the following statements about Huffman trees is/are true?
    \begin{choices}
        \choice Every internal node in a standard binary Huffman tree has exactly two children.
        \choice For $n$ input symbols, Huffman's algorithm performs exactly $n-1$ combine operations.
        \choice A Huffman tree with $n$ different characters has
        exactly $n$ leaves.
        \choice Every Huffman tree must be a complete binary tree.
    \end{choices}

\part[3] Which of the following statements about AVL trees is/are true?
    \begin{choices}
        \choice At every node of an AVL tree, the absolute difference between the heights of its left and right subtrees is at most $1$.
        \choice Every AVL tree is also a complete binary tree.
        \choice Searching for a key in an AVL tree takes $O(\log n)$ time in the worst case.
        \choice For any AVL tree, inserting one key into it may require rotations at every ancestor of the inserted node.
    \end{choices}

\end{parts}
